\originalchapter{相对论量子力学：从守恒流到Dirac氢原子}
本章对应8.322相对论合本全部34页。物理动机是非相对论$E=p^2/(2m)$不能适用于高能粒子；但用相对论色散替换Hamiltonian还会引入负能与多粒子问题。因此须把数学上可解的波方程、可解释的守恒流和完整量子场理论的边界区分。除明确恢复常量的公式，本章中间推导取$\hbar=c=1$，度规$\eta=\diag(1,-1,-1,-1)$。

\originalsection{Klein--Gordon方程及其解释边界}
四坐标$x^\mu=(ct,\vct x)$、四动量$p^\mu=(E/c,\vct p)$满足$p^2=m^2c^2$。Lorentz变换$\Lambda^T\eta\Lambda=\eta$，速度$x$向的惯性系变换给$E'=\gamma(E-vp_x)$、$p_x'=\gamma(p_x-vE/c^2)$。替换$E\to i\hbar\partial_t$、$\vct p\to-i\hbar\nabla$得到
\[
 (\Box+m^2)\phi=0,\qquad
 \Box=\partial_t^2-\nabla^2.
\]
复共轭方程相减得守恒流$j^\mu=i[\phi^*\partial^\mu\phi-(\partial^\mu\phi^*)\phi]/(2m)$。正频平面波$e^{-iEt+i\vct p\cdot\vct x}$给$j^0=(E/m)|\phi|^2$；负频给负号，且一般叠加也不保证点态密度正。这适宜解释为电荷流，不能同时当正概率密度和普通单粒子Schrödinger态。二阶时间方程需要初值$\phi,\dot\phi$，并非只给$\phi$就能演化。标量场量子化以正负频分别对应粒子、反粒子，而不是删掉一半解。

\originalsection{线性化与Clifford代数}
Dirac寻找$i\partial_t\psi=(\vct\alpha\cdot\vct p+\beta m)\psi$，平方要回到$E^2=p^2+m^2$。因此$\{\alpha_i,\alpha_j\}=2\delta_{ij}$、$\{\alpha_i,\beta\}=0$、$\beta^2=1$。四个反对易矩阵不能用二维Pauli矩阵满足：与三个Pauli全反对易的非零二维矩阵不存在。最小表示四维，可取
\[
 \beta=\begin{pmatrix}I&0\\0&-I\end{pmatrix},\quad
 \alpha_i=\begin{pmatrix}0&\sigma_i\\\sigma_i&0\end{pmatrix},\quad
 \gamma^0=\beta,\quad
 \gamma^i=\beta\alpha_i=\begin{pmatrix}0&\sigma_i\\-\sigma_i&0\end{pmatrix}.
\]
它们满足$\{\gamma^\mu,\gamma^\nu\}=2\eta^{\mu\nu}I$；方程$(i\gamma^\mu\partial_\mu-m)\psi=0$左乘共轭一阶因子，每分量满足KG。反之任意四分量KG解未必满足Dirac的一阶约束。定义$\bar\psi=\psi^\dagger\gamma^0$，乘方程与伴随方程得$\partial_\mu(\bar\psi\gamma^\mu\psi)=0$。密度$j^0=\psi^\dagger\psi\ge0$、流$\vct j=\psi^\dagger\vct\alpha\psi$；在外场仅讨论保持单粒子解释的弱场条件，强场负能问题仍需量子场。

四个$\gamma$各取或不取，按固定顺序相乘得到$1+4+6+4+1=16$个独立矩阵，PS12给迹正交证明。这也给Dirac矩阵表示的唯一性：它们张成$M_4(\mathbb C)$，两套满足相同Clifford关系的不可约表示诱导矩阵代数的自同构；用矩阵单位构造基变换，使该同构成为$A\mapsto SAS^{-1}$。相似矩阵可再按Hermiticity条件选酉规范，不把任意相似矩阵自动称为酉。

\originalsection{Lorentz协变与双线性量}
旋转生成元$\Sigma_i/2$，$\Sigma_i=\diag(\sigma_i,\sigma_i)$；$x$向被动boost的旋量变换为$S_x(\eta)=e^{-\eta\alpha_x/2}$，$\tanh\eta=v/c$。它不是酉矩阵，却满足$S^{-1}\gamma^\mu S=\Lambda^\mu{}_{\nu}\gamma^\nu$和$S^\dagger\gamma^0=\gamma^0S^{-1}$，故Dirac方程与四流协变。空间体积元和同时面也变换，不能仅由四分量欧氏模长不变来论证Lorentz协变。实验室旋转中的自旋归一与Lorentz变换的协变归一是不同目的。

令$\gamma^5=i\gamma^0\gamma^1\gamma^2\gamma^3$、$\sigma^{\mu\nu}=i[\gamma^\mu,\gamma^\nu]/2$。16个双线性按Lorentz性质分为标量$\bar\psi\psi$、赝标量$\bar\psi i\gamma^5\psi$、矢量$\bar\psi\gamma^\mu\psi$、轴矢量$\bar\psi\gamma^\mu\gamma^5\psi$和反对称张量$\bar\psi\sigma^{\mu\nu}\psi$。其协变性由上述相似变换逐项得出；宇称额外区分赝量。四分量旋量包含两个正能与两个负能自旋自由度，不能说成四个独立空间方向的经典向量。

Poincaré代数包含旋转$J_i$、boost $K_i$和平移$P^\mu$；在$\hbar=1$约定下
\[
 [J_i,J_j]=i\epsilon_{ijk}J_k,\quad[J_i,K_j]=i\epsilon_{ijk}K_k,\quad
 [K_i,K_j]=-i\epsilon_{ijk}J_k,\quad[P^\mu,P^\nu]=0.
\]
剩余对易子由$P$为四矢量决定。Casimir为$P^2=m^2$及$W^2=-m^2s(s+1)$，$W^\mu=\epsilon^{\mu\nu\rho\sigma}P_\nu J_{\rho\sigma}/2$。正能有质量粒子的静止四动量被旋转小群保持，因此自旋由$SU(2)$表示标记；Dirac粒子是$s=1/2$。此论证解释质量、自旋的分类，不宣称已证明全量子场的相互作用分类。无质量小群和helicity需另外处理，不把静止系构造用于光子。

\par
\ifdim\dimexpr\pagegoal-\pagetotal\relax<5\baselineskip\newpage\fi
\originalsection{自由解、局域化与负能}
平面波$\psi=u(p)e^{-iEt+i\vct p\cdot\vct x}$，写$u=(\varphi,\chi)^T$，方程给$\chi=(\sigma\cdot p)/(E+m)\,\varphi$。对$E=E_p=\sqrt{p^2+m^2}$可取
\[
 u_s(p)=\sqrt{\frac{E_p+m}{2E_p}}
 \begin{pmatrix}\xi_s\\\frac{\sigma\cdot p}{E_p+m}\xi_s\end{pmatrix},\qquad
 u_s^\dagger u_{s'}=\delta_{ss'},\quad\bar u_su_{s'}=\frac m{E_p}\delta_{ss'}.
\]
负能$E=-E_p$同样有两解，选下两分量为基即可。完整四分量初值展开不能只用固定动量两条正能旋量；删除负能时投影变成非局部算符。要局域到小于Compton尺度$\hbar/(mc)$，波包需$\Delta p\gtrsim mc$，相对论与粒子对自由度开始相关，不能从非相对论任意窄Gaussian的形式就保证可做单电子态。

势台阶上的Klein区是这种边界的直接表现。当$E-V<-mc^2$，右侧又出现传播负频解；选择右向群速度会使波数和电荷流为负，反射电荷系数可大于1。它不是概率不守恒，而是原“一个固定粒子”的解释不足；PS12将反射、透射流公式与这一限制一起给出。

\originalsection{电磁耦合、Pauli极限与氢细结构}
SI Dirac Hamiltonian为$H=c\vct\alpha\cdot\vct\pi+\beta mc^2+q\phi$，$\vct\pi=\vct p-q\vct A$。去掉静止相位$e^{-imc^2t/\hbar}$后写大分量$\varphi$与小分量$\chi$，若动能、场势及变化速率远小于$mc^2$，小分量方程给$\chi\simeq(\sigma\cdot\pi)\varphi/(2mc)$。利用
\[
 (\sigma\cdot\pi)^2=\pi^2-q\hbar\sigma\cdot\vct B
\]
得到$i\hbar\dot\varphi=[\pi^2/(2m)+q\phi-q\hbar\sigma\cdot\vct B/(2m)]\varphi$，$g=2$是Dirac树级结果。反常磁矩是量子场修正，未包含在这一消元近似中。更高阶慢场展开给质量--速度$-p^4/(8m^3c^2)$、自旋轨道$\hbar\sigma\cdot(\nabla V\times p)/(4m^2c^2)$与Darwin项$\hbar^2\nabla^2V/(8m^2c^2)$；本科氢细结构篇已逐项积分，不能只保留自旋轨道而称完整细结构。

中心势$V(r)$保$J=L+S$与$K=\beta(\sigma\cdot L+1)$。为避免与一些书的相反符号混用，定义$K\psi=-\kappa\psi$、$\kappa=\pm(j+1/2)$，自旋球谐$\Omega_{\kappa m}$满足$(\sigma\cdot L+1)\Omega_{\kappa m}=-\kappa\Omega_{\kappa m}$以及$\sigma\cdot\hat r\,\Omega_{\kappa m}=-\Omega_{-\kappa m}$。总自旋球谐由上一分编的$l\otimes1/2$系数构造。取
\[
 \psi=\frac1r\begin{pmatrix}F(r)\Omega_{\kappa m}\\iG(r)\Omega_{-\kappa m}\end{pmatrix},\quad
 F'+\frac\kappa rF=(E+m-V)G,\quad
 G'-\frac\kappa rG=(m-E+V)F.
\]
这由$\sigma\cdot p$在球谐上的作用直接推出，两分量有相反轨道宇称、相同$j,m$。标量$L^2,S^2$不是完整Dirac态分别的共同量子数。

点核Coulomb $V=-Z\alpha/r$在$Z\alpha<|\kappa|$时，原点级数$F,G\sim r^s$的二元行列式给$s=\sqrt{\kappa^2-(Z\alpha)^2}$。束缚大$r$取$e^{-\rho r}$、$\rho=\sqrt{m^2-E^2}$，写$F,G=r^se^{-\rho r}$乘幂级数；代入径向式消去两组系数，记$\zeta=Z\alpha$，具体两组递推为
\[
 \begin{aligned}
 (s+n+\kappa)f_n-\zeta g_n&=\rho f_{n-1}+(E+m)g_{n-1},\\
 \zeta f_n+(s+n-\kappa)g_n&=\rho g_{n-1}+(m-E)f_{n-1},
 \end{aligned}
\]
$n=0$右边取0，正是上述原点行列式；从$n=1$逐次确定系数。终止于$n_r$时下一阶要求$\rho f_{n_r}+(E+m)g_{n_r}=0$和$\rho g_{n_r}+(m-E)f_{n_r}=0$，再与递推相容，给$\zeta E/\rho=n_r+s$。于是恢复常量，
\[
 E_{n\kappa}=mc^2\left[1+\frac{(Z\alpha)^2}
 {(n_r+\sqrt{\kappa^2-(Z\alpha)^2})^2}\right]^{-1/2},\qquad
 n=n_r+|\kappa|.
\]
$n_r=0$仅允许$\kappa<0$的相应正常支，不能凭能量公式创造$n=1,p_{1/2}$态。级数在无穷不终止会引入增长支，原点条件与无穷条件共同量子化；这说明能谱不是从色散式猜出的。弱$Z\alpha$展开为
\[
 E=mc^2-\frac{mc^2(Z\alpha)^2}{2n^2}
 -\frac{mc^2(Z\alpha)^4}{2n^4}
 \left[\frac n{j+1/2}-\frac34\right]+O((Z\alpha)^6).
\]
它汇合非相对论束缚能与全部细结构，仍不含核尺寸、反冲、超精细或Lamb位移。点核大$Z$时定义域和物理有限核都需重新审查，不将公式无条件外推。

\originalsection{原扫描逐段定位}
\begin{longtable}{p{.19\textwidth}p{.71\textwidth}}
\toprule 原页 & 本章对应\\\midrule\endhead
1--5 & Lorentz四矢量、KG、负频与守恒流。\\
6--11 & Dirac线性化、矩阵、正密度、16矩阵与表示等价。\\
12--19 & Poincare与质量/自旋、协变性、旋转/boost、双线性量。\\
20--25 & 自由正/负能、归一、局域化及Klein区。\\
26--29 & 负能解释边界、电磁耦合及Pauli极限。\\
30--34 & 总角动量、$\kappa$、球谐、径向Dirac和Coulomb能谱。\\
\bottomrule\end{longtable}
原讲义公式的单位、$\kappa$和boost约定在此固定后重推，不把来源中记号差异偷偷合并。完整作业题面及具体可计算例子接于下面两章。
